\begin{thebibliography}{99}\small\setlength{\itemsep}{0pt}
\bibitem{tutorial} M. E. P. Davies, S. B\"ock, and M. Fuentes, ``Tempo, beat and downbeat estimation,'' Tutorial at ISMIR, 2021. [Online]. Available: https://tempobeatdownbeat.github.io/tutorial
\bibitem{bello} J. P. Bello, L. Daudet, S. Abdallah, C. Duxbury, M. Davies, and M. Sandler, ``A tutorial on onset detection in music signals,'' \emph{IEEE Trans. Speech Audio Process.}, vol. 13, no. 5, pp. 1035--1047, 2005.
\bibitem{gouyon} F. Gouyon, A. Klapuri, S. Dixon, M. Alonso, G. Tzanetakis, C. Uhle, and P. Cano, ``An experimental comparison of audio tempo induction algorithms,'' \emph{IEEE Trans. Audio, Speech, Lang. Process.}, vol. 14, no. 5, pp. 1832--1844, 2006.
\bibitem{percival} G. Percival and G. Tzanetakis, ``Streamlined tempo estimation based on autocorrelation and cross-correlation with pulses,'' \emph{IEEE/ACM Trans. Audio, Speech, Lang. Process.}, vol. 22, no. 12, pp. 1765--1776, 2014.
\bibitem{ellis} D. P. W. Ellis, ``Beat tracking by dynamic programming,'' \emph{J. New Music Res.}, vol. 36, no. 1, pp. 51--60, 2007.
\bibitem{librosa} B. McFee, C. Raffel, D. Liang, D. P. W. Ellis, M. McVicar, E. Battenberg, and O. Nieto, ``librosa: Audio and music signal analysis in Python,'' in \emph{Proc. Python in Science Conf. (SciPy)}, pp. 18--25, 2015.
\bibitem{plp} P. Grosche and M. M\"uller, ``Extracting predominant local pulse information from music recordings,'' \emph{IEEE Trans. Audio, Speech, Lang. Process.}, vol. 19, no. 6, pp. 1688--1701, 2011.
\bibitem{downbeat} S. B\"ock, F. Krebs, and G. Widmer, ``Joint beat and downbeat tracking with recurrent neural networks,'' in \emph{Proc. ISMIR}, 2016.
\bibitem{madmom} S. B\"ock, F. Korzeniowski, J. Schl\"uter, F. Krebs, and G. Widmer, ``madmom: A new Python audio and music signal processing library,'' in \emph{Proc. ACM Multimedia}, pp. 1174--1178, 2016.
\bibitem{krebsbar} F. Krebs, S. B\"ock, M. Dorfer, and G. Widmer, ``Downbeat tracking using beat-synchronous features and recurrent neural networks,'' in \emph{Proc. ISMIR}, 2016.
\bibitem{tcn} M. E. P. Davies and S. B\"ock, ``Temporal convolutional networks for musical audio beat tracking,'' in \emph{Proc. EUSIPCO}, 2019.
\bibitem{heydari} M. Heydari, M. McCallum, A. Ehmann, and Z. Duan, ``A novel 1D state space for efficient music rhythmic analysis,'' in \emph{Proc. ICASSP}, 2022.
\bibitem{degradation} M. Mauch and S. Ewert, ``The Audio Degradation Toolbox and its application to robustness evaluation,'' in \emph{Proc. ISMIR}, pp. 83--88, 2013.
\bibitem{daviseval} M. E. P. Davies, N. Degara, and M. D. Plumbley, ``Evaluation methods for musical audio beat tracking algorithms,'' Queen Mary Univ. of London, Tech. Rep. C4DM-TR-09-06, 2009.
\bibitem{mireval} C. Raffel, B. McFee, E. J. Humphrey, J. Salamon, O. Nieto, D. Liang, and D. P. W. Ellis, ``mir\_eval: A transparent implementation of common MIR metrics,'' in \emph{Proc. ISMIR}, 2014.
\bibitem{gtzan} G. Tzanetakis and P. Cook, ``Musical genre classification of audio signals,'' \emph{IEEE Trans. Speech Audio Process.}, vol. 10, no. 5, pp. 293--302, 2002.
\bibitem{sturm} B. L. Sturm, ``The GTZAN dataset: Its contents, its faults, their effects on evaluation, and its future use,'' arXiv:1306.1461, 2013.
\bibitem{kereliuk} C. Kereliuk, B. L. Sturm, and J. Larsen, ``Deep learning and music adversaries,'' \emph{IEEE Trans. Multimedia}, vol. 17, pp. 2059--2072, 2015.
\bibitem{marchand} U. Marchand, Q. Fresnel, and G. Peeters, ``GTZAN-Rhythm: Extending the GTZAN test-set with beat, downbeat and swing annotations,'' in \emph{Late-Breaking Demo, ISMIR}, 2015.
\bibitem{giantsteps} P. Knees, \'A. Faraldo, P. Herrera, R. Vogl, S. B\"ock, F. H\"orschl\"ager, and M. Le Goff, ``Two data sets for tempo estimation and key detection in electronic dance music annotated from user corrections,'' in \emph{Proc. ISMIR}, 2015.
\end{thebibliography}